\section[Cap\'{i}tulo 1: Ciberseguridad en la Educaci\'{o}n Superior]{\centering Cap\'{i}tulo 1: Ciberseguridad en la Educaci\'{o}n Superior}

La ciberseguridad es el punto de partida del proyecto porque define el tipo de problemas que el videojuego busca representar. En una carrera inform\'{a}tica, este campo no deber\'{i}a tratarse solamente como teor\'{i}a, sino como una competencia pr\'{a}ctica para proteger informaci\'{o}n, sistemas y servicios frente a riesgos digitales (NIST, 2024).

Este cap\'{i}tulo re\'{u}ne conceptos base sobre seguridad de la informaci\'{o}n, amenazas, vulnerabilidades y factor humano. Estos elementos permiten justificar por qu\'{e} un simulador puede apoyar la formaci\'{o}n de estudiantes que necesitan analizar casos, reconocer se\~{n}ales de riesgo y decidir respuestas apropiadas (ISO/IEC, 2022).

\subsection{Concepto de ciberseguridad}

La ciberseguridad comprende pr\'{a}cticas, procesos y controles orientados a proteger sistemas, redes, dispositivos e informaci\'{o}n frente a accesos no autorizados, alteraciones, interrupciones o ataques. Sus principios m\'{a}s conocidos son confidencialidad, integridad y disponibilidad, porque resumen qu\'{e} se busca proteger en un sistema digital (Stallings \& Brown, 2018).

\begin{center}
\begin{tabular}{|p{0.24\textwidth}|p{0.58\textwidth}|}
\hline
\textbf{Principio} & \textbf{Relaci\'{o}n con el proyecto} \\
\hline
Confidencialidad & Evitar que credenciales, datos personales o informaci\'{o}n institucional sean expuestos por enga\~{n}os o accesos indebidos (NIST, 2024). \\
\hline
Integridad & Reconocer alteraciones en mensajes, archivos, registros o procesos que podr\'{i}an indicar actividad maliciosa (Stallings \& Brown, 2018). \\
\hline
Disponibilidad & Comprender que una mala decisi\'{o}n puede afectar servicios, equipos o continuidad de trabajo dentro de un escenario simulado (ISO/IEC, 2022). \\
\hline
\end{tabular}
\end{center}

En el videojuego, estos principios se traducen en situaciones concretas: un correo falso puede comprometer credenciales, un archivo malicioso puede modificar informaci\'{o}n y un equipo infectado puede dejar de operar. De esta forma, la teor\'{i}a se relaciona con decisiones visibles dentro de un escenario de juego (NIST, 2024).

\subsection{Amenazas y vulnerabilidades}

Una amenaza digital es una acci\'{o}n o evento capaz de causar da\~{n}o, mientras que una vulnerabilidad es una debilidad que puede ser aprovechada. Esta diferencia es importante porque permite analizar un incidente no solo desde el ataque, sino tambi\'{e}n desde las condiciones que facilitaron que ocurra (ISO/IEC, 2022).

Para el proyecto se consideran amenazas iniciales que pueden representarse con claridad en un videojuego educativo y que aparecen de forma recurrente en reportes de incidentes de seguridad (Verizon, 2024):

\begin{itemize}
    \item \textit{Phishing}: mensajes que imitan comunicaciones leg\'{i}timas para robar informaci\'{o}n o credenciales (Verizon, 2024).
    \item Malware: software malicioso que altera, bloquea, roba o destruye informaci\'{o}n (Stallings \& Brown, 2018).
    \item Ingenier\'{i}a social: manipulaci\'{o}n de personas para obtener accesos, c\'{o}digos o datos sensibles (Verizon, 2024).
    \item Ransomware: ataque que cifra informaci\'{o}n y exige pago para recuperarla (Verizon, 2024).
\end{itemize}

Las vulnerabilidades no siempre dependen de fallas t\'{e}cnicas. Tambi\'{e}n pueden surgir por contrase\~{n}as d\'{e}biles, sistemas sin actualizar, permisos excesivos o falta de verificaci\'{o}n de enlaces. Por eso, el videojuego debe presentar pistas que obliguen al estudiante a mirar detalles antes de decidir (Stallings \& Brown, 2018).

\subsection{Factor humano en los incidentes}

El factor humano tiene un peso importante en los incidentes de seguridad, especialmente cuando existen credenciales robadas, errores de configuraci\'{o}n o respuestas apresuradas ante mensajes enga\~{n}osos. Esta realidad hace necesario trabajar la formaci\'{o}n desde escenarios donde el estudiante pueda equivocarse y revisar consecuencias sin afectar sistemas reales (Verizon, 2024).

En el proyecto, este aspecto se refleja en correos sospechosos, llamadas de supuesto soporte, enlaces con dominios alterados y decisiones de contenci\'{o}n. La intenci\'{o}n es que el estudiante practique una forma de razonamiento profesional: observar, comparar, sospechar con fundamento y actuar seg\'{u}n el nivel de riesgo (NIST, 2024).

\subsection{S\'{i}ntesis del cap\'{i}tulo}

La ciberseguridad aporta el contenido central que el videojuego debe representar. Los conceptos de amenaza, vulnerabilidad, riesgo y factor humano permiten dise\~{n}ar casos donde el estudiante no solo reconoce definiciones, sino que aplica criterios de seguridad en situaciones concretas.
